\section{Arithmetic ledgers and short-interval transfer}
\label{sec:arithmetic-transfer}

The geometric lift changes the spectral fibres but not the product
centers to which the short-interval theorem is applied.  For $r\geq1$
put
\begin{align}
 A_D(r)
 &:={}|W(r/X)|^2
 \sum_{\substack{n\mid r\\n>D}}u(n)^2,
 \label{eq:product-diagonal-weight}\\
 B^{\rad}_{D,t,s}(r)
 &:={}W(r/X)
 \sum_{\substack{n\mid r\\n>D}}
 u(n)(n^2/r)^{it}
 \gcd(n,r/n)^{is}.
 \label{eq:radial-signed-product-weight}
\end{align}
The first weight is unchanged from TPC-12.  The second is the coefficient
obtained by grouping \eqref{eq:intro-centered-radial-transform} by
$r=mn$.

\begin{lemma}[Product ledgers]
\label{lem:radial-product-ledgers}
Uniformly in $D\geq1$ and $t,s\in\R$,
\begin{align}
 A_D(r)&\ll_W\one_{[\alpha X,\beta X]}(r)
 \bigl((\log(2r))^2+\tau(r)\bigr),
 \label{eq:A-pointwise}\\
 |B^{\rad}_{D,t,s}(r)|
 &\ll_W\one_{[\alpha X,\beta X]}(r)
 \bigl(\log(2r)+\tau(r)\bigr),
 \label{eq:B-radial-pointwise}
\end{align}
and
\begin{align}
 \sum_rA_D(r)^2&\ll_WX(\log(2X))^4,
 \label{eq:A-second-moment}\\
 \sum_r|B^{\rad}_{D,t,s}(r)|^2
 &\ll_WX(\log(2X))^3.
 \label{eq:B-radial-second-moment}
\end{align}
\end{lemma}

\begin{proof}
Since $(\Lambda-1)^2\leq2(1+\Lambda^2)$,
\[
 \sum_{d\mid r}u(d)^2
 \leq2\tau(r)+2\sum_{p^j\mid r}(\log p)^2
 \leq2\tau(r)+2(\log r)^2.
\]
Also, both phases in \eqref{eq:radial-signed-product-weight} have modulus
one, so
\[
 |B^{\rad}_{D,t,s}(r)|
 \leq\|W\|_\infty\sum_{d\mid r}|u(d)|
 \leq\|W\|_\infty\bigl(\log r+\tau(r)\bigr).
\]
The second moments follow from
$\sum_{r\leq y}\tau(r)^2\ll y(\log(2y))^3$.  This is the same proof as
the product ledger in TPC-12; the new GCD phase costs nothing.
\end{proof}

Define the total source diagonal
\begin{equation}
 \cS_D(X):=\sum_rA_D(r)
 =\sum_{\substack{n>D\\m\geq1}}
 u(n)^2|W(mn/X)|^2.
 \label{eq:source-diagonal-definition}
\end{equation}

\begin{lemma}[Source diagonal]
\label{lem:source-diagonal}
For every fixed $\eta>0$, uniformly for
$1\leq D\leq X^{1-\eta}$,
\begin{equation}
 \cS_D(X)
 =\frac12XI_2(W)
 \bigl((\log X)^2-(\log D)^2\bigr)
 +O_W(X\log(2X)).
 \label{eq:source-diagonal}
\end{equation}
\end{lemma}

\begin{proof}
Euler summation gives
\[
 \sum_{m\geq1}|W(mn/X)|^2
 =\frac XnI_2(W)+O_W(1).
\]
The prime number theorem and partial summation give
\[
 \sum_{n\leq y}\frac{(\Lambda(n)-1)^2}{n}
 =\frac12(\log y)^2-\log y+O(1),
\]
while $\sum_{n\leq y}u(n)^2\ll y\log(2y)$.  Subtracting the range
$n\leq D$ proves the assertion; see also
\citet[Chapter~5]{IwaniecKowalski2004} and TPC-12
\citep{WangTPC12}.
\end{proof}

The following transferred moment is the deep arithmetic input.  It is
quoted in the form proved as Proposition~4.2 of TPC-12.  Its proof starts
from Corollary~1.4 of arXiv:2405.20552v2, the Guth--Maynard prime number
theorem for almost all intervals of length
$X^{2/15+\eps}$ \citep{GuthMaynard2026}.  The exceptional set has size
$O(X\exp\{-c(\log X)^{1/4}\})$ for each discretized interval length;
TPC-12 takes a controlled union over the lengths needed to smooth $G$.
Cauchy--Schwarz and \eqref{eq:A-second-moment} then show that the
exceptional centers carry negligible $A_D$-mass.
For $H\leq X^{0.98}$ this is the almost-all argument.  For
$H>X^{0.98}$, TPC-12 partitions the full shift interval into equal
blocks of lengths in $[X^{0.97},2X^{0.97}]$ and applies the uniform
$17/30$ short-interval theorem on every block; there is no exceptional
center in that range.  These two cases cover the full upper range
$H\leq X$ stated below.

\begin{proposition}[Weighted target-square transfer]
\label{prop:weighted-target-square}
Let $G\in C_c^1((0,\infty))$ be nonnegative and nonzero.  Fix
$0<\eps<13/15$ and $0<\eta<1$.  Uniformly for
\eqref{eq:intro-parameter-range},
\begin{equation}
 \begin{aligned}
 &\sum_rA_D(r)
 \sum_hG(h/H)(\Lambda(r+h)-1)^2\\
 &\qquad={}
 HI_0(G)\log X\,\cS_D(X)
 +O_{\eps,W,G}\bigl(HX(\log X)^2\bigr).
 \end{aligned}
 \label{eq:weighted-centered-target-square}
\end{equation}
The same formula holds with the target square replaced by
$\Lambda(r+h)^2$.
\end{proposition}

\begin{proof}
This is Proposition~4.2 of \citet{WangTPC12}.  Outside the transferred
exceptional set, smooth
partial summation gives
\[
 \sum_hG(h/H)(\Lambda(r+h)-1)^2
 =HI_0(G)\log X+O(H).
\]
The $O(H)$ term sums to $O(H\cS_D)=O(HX(\log X)^2)$.  On the exceptional
set, the pointwise $O(H(\log X)^2)$ bound and the arbitrarily large
logarithmic saving in its $A_D$-mass give a smaller contribution.  The
raw square is identical.
\end{proof}

The signed transfer also survives the new phase.

\begin{theorem}[Radially completed short-shift first moment]
\label{thm:radial-first-moment}
Under \eqref{eq:intro-parameter-range}, for every $A>0$, uniformly in
$t,s\in\R$,
\begin{equation}
 \frac1H\sum_hG(h/H)\cT^{\rad}_{h,D}(X;t,s)
 =I_0(G)\cM^{\rad}_D(X;t,s)
 +O_{A,\eps,W,G}\bigl(X(\log X)^{-A}\bigr).
 \label{eq:radial-first-moment}
\end{equation}
\end{theorem}

\begin{proof}
Grouping by $r=mn$ gives
\[
 \frac1H\sum_hG(h/H)\cR^{\rad}_{h,D}(X;t,s)
 =\sum_rB^{\rad}_{D,t,s}(r)
 \left\{\frac1H\sum_hG(h/H)(\Lambda(r+h)-1)\right\}.
\]
On regular centers the braces save every fixed power of $\log X$.
Equation \eqref{eq:B-radial-pointwise} and the mean order of $\tau$
give $\sum_r|B^{\rad}_{D,t,s}(r)|\ll_WX\log(2X)$.  On the exceptional
set, Cauchy--Schwarz, \eqref{eq:B-radial-second-moment}, and its
stretched-exponential cardinality again save every logarithmic power.
All bounds are uniform in both frequency variables because the added
phase has modulus one.  Finally the Riemann-sum normalization of $G$
produces the displayed model term.
\end{proof}
